\documentclass[10pt,a4paper]{amsart}
\usepackage[margin=19mm]{geometry}
\usepackage{amsmath,amssymb,mathtools}
\usepackage[scr=rsfs,cal=euler]{mathalfa}
\usepackage{xfrac}
\usepackage[unicode,hidelinks,bookmarks=false]{hyperref}
\newcommand{\ie}{i.e.\ }
\newcommand{\cf}{cf.\ }
\newcommand{\resp}{respectively\ }
\newcommand{\Subsection}{\subsection}
\providecommand{\nobreakdash}{\nobreak\hbox{-}}
\setlength{\emergencystretch}{2em}
\multlinegap0pt
\usepackage{bm}
\usepackage{bbm}
\usepackage{dsfont}
\usepackage{xspace}
\usepackage{cancel}
\usepackage{mathtools}
\usepackage{amsthm}
\makeatletter
\@ifundefined{thm}{\newtheorem{thm}{Thm}}{}
\@ifundefined{lemma}{\newtheorem{lemma}[thm]{Lemma}}{}
\makeatother
\newcommand{\BCH}{\operatorname{BCH}}
\title[A Quaternion--BCH Framework for the Local Accuracy of SIDER ]{A Quaternion--BCH Framework for the Local Accuracy of SIDER Interpolation}
\author{Shingyu Leung}
\date{Source: arXiv:2606.12833v1 (CC BY 4.0)}
\begin{document}
\maketitle

\noindent\emph{Source and licence.} A Quaternion--BCH Framework for the Local Accuracy of SIDER Interpolation. Version arXiv:2606.12833v1 (CC BY 4.0), distributed under the Creative Commons Attribution 4.0 International licence, as recorded in the arXiv licence metadata for this exact version and shown on the abstract page https://arxiv.org/abs/2606.12833v1. Full rights notice: licenses/arxiv-2606.12833-CC-BY-4.0.txt.\par\smallskip

\noindent\textit{Editorial scope.} Counted unit: the Lemma whose source label is \texttt{lem:filtered\_bch\_closure} in the article (source0.tex lines 1171--1204, source label \texttt{lem:filtered\_bch\_closure}), delivered together with the article's own complete original proof of it (source0.tex lines 1206--1377, one \texttt{proof} environment of 172 lines). No mathematics is written, repaired or re-derived: statement and proof are the article's own text line for line. Required context retained uncounted: none. Every definition, display and label of those lines is retained unchanged. Omitted from this package and not redistributed: 1--1170, 1205--1205, 1378--1859. Nothing inside a retained excerpt is omitted or altered: every retained body is delivered exactly as the article's lines are written, so no omission receipt of any kind accompanies this package. Citations: the retained excerpts carry no citation command at all, so no bibliography entry of the article is transcribed and no reference is redistributed. Reference adaptation: none was needed and none was made; every reference inside the delivered unit resolves inside it, so no unresolved reference or citation is produced. Printed slips kept literal: no printed word, symbol, hypothesis or number of the counted statement or of its proof was altered. Layout and class: the article's own class is \texttt{article} with packages including grffile, latexsym, amssymb, enumerate, amsmath, epsfig, amsthm, geometry, setspace, color, graphics, graphicx, algorithm2e, empheq; the wrapper is the lane's pinned \texttt{amsart} editorial wrapper at 10pt on A4, which restates the theorem-like environments the retained text needs and adds the article's own macro definitions that the retained text needs (\texttt{\textbackslash BCH}), with extra wrapper packages (bm, bbm, dsfont, xspace, cancel, mathtools). The delivered numbering is the unit's own. Assets: no retained span contains a figure environment, a graphics-inclusion command, a PDF-page inclusion, a table or a TikZ picture; no image file is copied into this package, the package declares no asset dependency and no image file is redistributed with it. Licence: version arXiv:2606.12833v1 of this article is distributed under the Creative Commons Attribution 4.0 International licence, as recorded in the arXiv licence metadata for this exact version and shown on the abstract page https://arxiv.org/abs/2606.12833v1. A full rights notice is delivered at licenses/arxiv-2606.12833-CC-BY-4.0.txt. Characters: every retained excerpt is pure ASCII, so the delivered engine, XeLaTeX with its default Latin Modern text font, reports no missing character.
\begin{lemma}[Filtered closure of the nonlinear BCH remainder]
\label{lem:filtered_bch_closure}
Let \(k\geq3\).  Suppose
\[
        X(\xi,h)
        =
        h^k\Pi_{k-1}(\xi)C(\theta,\xi,h),
\]
and
\[
        Y(\xi,h)
        =
        h^k\Pi_{k-1}(\xi-1)C(\theta,\xi-1,h),
\]
where \(C\) is degree-filtered.  Set
\[
        \lambda=\frac{\xi}{k}
\]
and define the nonlinear SLERP remainder by
\[
        R_\lambda(X,Y)
        =
        \mathcal S_\lambda(X,Y)
        -
        \bigl((1-\lambda)X+\lambda Y\bigr).
\]
Then
\[
        R_{\xi/k}(X,Y)
        =
        h^{k+1}\Pi_k(\xi)\mathcal R_k(\theta,\xi,h),
\]
where \(\mathcal R_k\) is degree-filtered.
\end{lemma}

\begin{proof}
We use the diagonal variable
\[
        \Delta=Y-X.
\]
Thus \(Y=X+\Delta\), and
\[
        R_\lambda(X,Y)
        =
        \mathcal S_\lambda(X,X+\Delta)-X-\lambda\Delta .
\]
Define
\[
        \mathcal N_\lambda(X,\Delta)
        =
        \mathcal S_\lambda(X,X+\Delta)-X-\lambda\Delta .
\]
Then
\[
        R_\lambda(X,Y)=\mathcal N_\lambda(X,Y-X).
\]

We first describe the formal structure of \(\mathcal N_\lambda\).  Let
\[
        Z(X,\Delta)=\BCH(-X,X+\Delta).
\]
Since
\[
        Z(X,0)=\BCH(-X,X)=0,
\]
every term in the formal expansion of \(Z(X,\Delta)\) contains at least one
factor of \(\Delta\).  Moreover,
\[
        \mathcal S_\lambda(X,X+\Delta)
        =
        \BCH\bigl(X,\lambda Z(X,\Delta)\bigr).
\]
The constant term in \(\Delta\) is \(X\).  The term linear in \(\Delta\) is
\(\lambda\Delta\).  To see this, observe that
\[
        \BCH\bigl(X,Z(X,\varepsilon\Delta)\bigr)
        =
        X+\varepsilon\Delta
\]
for \(\varepsilon\) sufficiently small, because
\[
        \exp X\exp Z(X,\varepsilon\Delta)
        =
        \exp(X+\varepsilon\Delta).
\]
Differentiating this identity at \(\varepsilon=0\) gives
\[
        D_2\BCH(X,0)\,Z_\varepsilon(X,0)=\Delta .
\]
Therefore
\[
        \frac{d}{d\varepsilon}
        \BCH\bigl(X,\lambda Z(X,\varepsilon\Delta)\bigr)
        \bigg|_{\varepsilon=0}
        =
        \lambda\Delta .
\]
After subtracting \(X+\lambda\Delta\), every term in
\(\mathcal N_\lambda(X,\Delta)\) contains at least two factors of
\(\Delta\).

We also need to track the dependence on \(\lambda\).  Each occurrence of
\(\lambda\) in
\[
        \BCH\bigl(X,\lambda Z(X,\Delta)\bigr)
\]
comes from an occurrence of the second BCH argument
\(\lambda Z(X,\Delta)\).  Since each factor \(Z(X,\Delta)\) contains at least
one factor of \(\Delta\), the degree in \(\lambda\) of any Lie monomial in
\(\mathcal N_\lambda\) is no larger than the number of its \(\Delta\)-factors.

We now apply this structure to the adjacent SIDER errors.  Since \(C\) is
degree-filtered, the coefficient of \(h^{k+r}\) in
\[
        X=h^k\Pi_{k-1}(\xi)C(\theta,\xi,h)
\]
has degree at most \(k+r\) in \(\xi\).  On the other hand,
\[
\begin{aligned}
        \Delta
        &=
        Y-X  \\
        &=
        h^k
        \left[
        \Pi_{k-1}(\xi-1)C(\theta,\xi-1,h)
        -
        \Pi_{k-1}(\xi)C(\theta,\xi,h)
        \right].
\end{aligned}
\]
The coefficient of \(h^{k+r}\) in \(\Delta\) has degree at most \(k+r-1\).
Indeed, \(\Pi_{k-1}(\xi-1)\) and \(\Pi_{k-1}(\xi)\) are monic polynomials of
degree \(k\), while \(C_r(\theta,\xi-1)\) and \(C_r(\theta,\xi)\) have the
same leading coefficient.  Hence the highest-degree terms cancel in the
difference.

Consider any Lie monomial appearing in
\[
        \mathcal N_{\xi/k}(X,\Delta).
\]
Suppose it contributes to the coefficient of \(h^s\), and suppose it contains
\(b\) factors of \(\Delta\).  Before including the powers of
\(\lambda=\xi/k\), the degree in \(\xi\) is at most
\[
        s-b,
\]
because each \(\Delta\)-factor gives one degree saving.  The degree in
\(\lambda\) is at most \(b\), so substituting \(\lambda=\xi/k\) raises the
degree in \(\xi\) by at most \(b\).  Therefore the total degree of the
coefficient of \(h^s\) is at most
\[
        (s-b)+b=s.
\]
Thus every coefficient of \(R_{\xi/k}(X,Y)\) has polynomial degree at most its
power of \(h\).

It remains to extract the nodal factor.  The nonlinear remainder vanishes at
all nodes
\[
        \xi=0,1,\ldots,k.
\]
At \(\xi=0\), we have \(\lambda=0\), so
\[
        \mathcal S_0(X,Y)=X
\]
and hence \(R_\lambda(X,Y)=0\).  At \(\xi=k\), we have \(\lambda=1\), so
\[
        \mathcal S_1(X,Y)=Y
\]
and again \(R_\lambda(X,Y)=0\).  For the interior nodes
\(\xi=1,\ldots,k-1\), both nodal factors
\[
        \Pi_{k-1}(\xi)
        \quad\text{and}\quad
        \Pi_{k-1}(\xi-1)
\]
vanish.  Hence \(X=Y=0\), and the nonlinear remainder also vanishes.

Therefore each coefficient in the formal expansion of \(R_{\xi/k}(X,Y)\) is
divisible by
\[
        \Pi_k(\xi)=\prod_{m=0}^{k}(\xi-m).
\]
Since the coefficient of \(h^s\) has degree at most \(s\), division by the
degree-\((k+1)\) polynomial \(\Pi_k\) gives a quotient coefficient of degree
at most
\[
        s-k-1.
\]

Finally, the nonlinear BCH remainder starts at cubic order in its inputs.
Since \(X,Y=O(h^k)\), we have
\[
        R_{\xi/k}(X,Y)=O(h^{3k}).
\]
In particular, for \(k\geq3\), this is of order at least \(h^{k+1}\).  Thus we
may write
\[
        R_{\xi/k}(X,Y)
        =
        h^{k+1}\Pi_k(\xi)\mathcal R_k(\theta,\xi,h).
\]
The coefficient of \(h^r\) in \(\mathcal R_k\) comes from the coefficient of
\(h^{r+k+1}\) in \(R_{\xi/k}\), and therefore has degree at most \(r\).  Hence
\(\mathcal R_k\) is degree-filtered.
\end{proof}

\end{document}
